\section{Cusp growth, component volume and the degree of the image}
\label{sec:component-volume}
We give the two inequalities behind \eqref{eq:rankshare-vol}. Let
\(m=p+q\), let \(\Phi\) be the selected set of complex embeddings of the
CM field, and set \(D=|\Phi|\binom m2\). For a weight \(w\) set
\[
 V(w)=\prod_{\tau\in\Phi}\prod_{0\le i< j<m}(w_{\tau i}-w_{\tau j}),
 \qquad \eta_{\tau i}=m-1-i .
\]
Only differences of \(\eta\) matter below. On a good deep component \(Z\),
let \(S_Z(rw)\) be the cusp space, taken at scalar twist zero, and put
\(g_Z(w)=D!\limsup_{r\to\infty}\dim S_Z(rw)/r^D\). The source's
extended nonnegative growth agrees with this expression once finite
dimensionality holds. All weights used in the volume comparisons are
strictly decreasing at each embedding.

\paragraph{Nefness and restriction to the boundary.}
On the smooth projective flag compactification let \(L_w\) be the weight
bundle. For a weakly decreasing integer vector \(w_0,\ldots,w_{m-1}\)
at one embedding put \(g_i=w_i-w_{i+1}\) for \(0\le i<m-1\).
Write \(h_i\) for the vector with ones in positions \(0,\ldots,i\) and
zero elsewhere, \(t_i=h_i-\mathbf 1\), and
\(c=w_{m-1}+\sum_{i\ge p}g_i\), where the sum ends at \(m-2\). Then
\[
 w=\sum_{i<p}g_i h_i+\sum_{p\le i<m-1}g_i t_i+c\mathbf 1
\]
with indices in both sums restricted to \(0\le i<m-1\).
Indeed the last coordinate is \(w_{m-1}\) and taking successive differences
recovers every \(g_i\). The source's head and tail bundle estimates give
nefness of \(L_{h_i}\) for \(i<p\) and \(L_{t_i}\) for \(i\ge p\).
Both determinant weights \(\mathbf 1,-\mathbf 1\) give nef bundles.
The weight-to-Picard-class map is additive. On every integral curve all
degrees in this decomposition are nonnegative, using \(|c|\) copies
of the appropriate determinant bundle. Summing over embeddings proves
that \(L_w\) is nef. The head, tail and determinant estimates use the
semiabelian and torus bundles on the compactification.

Let \(B\) be the reduced boundary. The extension and boundary-vanishing
theorems identify the cusp conditions in \(H^0(L_{rw})\). The kernel of
restriction to \(B\) lies in the cusp space, and every cusp section lies
in \(H^0(L_{rw})\). Thus
\[
 h^0(X,L_{rw})-h^0(B,L_{rw}|_B)
 \ \le\ \dim S_Z(rw)\ \le\ h^0(X,L_{rw}) .
\]
Here \(L_{rw}\simeq L_w^{\otimes r}\). The interior is dense, so
\(\dim B<D\). Nef asymptotic Riemann--Roch gives
\(h^0(X,L_w^{\otimes r})/r^D\to L_w^D/D!\), and, on the reduced
boundary, \(h^0(B,L_w^{\otimes r}|_B)/r^D\to0\).
The squeeze thus proves the existence of the cusp growth limit and
\begin{equation}\label{eq:volume-nef-growth}g_Z(w)=L_w^D.\end{equation}

\paragraph{Normalization of the flag degree.}
Use the compactification with toric boundary charts and the source's
Mumford proportionality formula. It identifies \(L_w^D\) with the
covolume \(c_Z\) of the uniformizing group, in the chosen measure,
times the degree polynomial of the compact dual flag. In the source's
alternating coefficient formula, expand the \(D\)-th power of the weight
linear form. For a single embedding its factor, after extracting the common
multinomial numerator \(D!\) across all embeddings, is
\[
 \sum_{\pi\in\mathfrak S_m}\operatorname{sgn}(\pi)
 \prod_{i=0}^{m-1}\frac{w_{\tau i}^{\,m-1-\pi(i)}}{(m-1-\pi(i))!}
 =\frac{\prod_{i<j}(w_{\tau i}-w_{\tau j})}
             {\prod_{j=0}^{m-1}j!}.
\]
The sum is a determinant with descending powers. Reversing its columns
and applying the Vandermonde determinant produces the displayed sign.
Since \(\prod_jj!=\prod_{i<j}(j-i)\), multiplying the embedding factors
gives
\begin{equation}\label{eq:volume-flag}
 g_Z(w)=c_Z D!\,\frac{V(w)}{V(\eta)} .
\end{equation}
The constant \(c_ZD!/V(\eta)\) is strictly positive. To compare two
components one uses the equal-covolume result at the actual deep level.
The good local level at each rational prime is inherited from the
original two-prime construction. Complex uniformization supplies
full-level self-dual data for each component with the same covolume as
that component. Their complex types and signatures agree. Transport to
one Hermitian space preserves the covolume; the self-dual lattice
comparison there identifies the two volumes. In particular the
right side of \eqref{eq:volume-flag} is independent of \(Z\) for fixed
\(w\). This use of the self-dual lattice theorem is separate from the
calculation of the flag degree.

\paragraph{The two inequalities for the image degree.}
Set \(u'=u'(u)\) equal to the lower weight of the window. Write
\(\delta_Z\) for the degree of the birational image of the complete cusp
system at weight \(u\). Then
\begin{equation}\label{eq:volume-sandwich}
 g_{Z'}(u')\ \le\ \delta_Z\ \le\ g_{Z'}(u)
 \qquad\text{for every pair }Z,Z'.
\end{equation}
We prove this first with \(Z'=Z\). Use the function field \(\Omega\)
of the generic component with the source's evaluation trivialization.
Write \(W\) for the image of \(S_Z(u)\) there and form the standard
graded domain \(R=\bigoplus_{j\ge0}W^j X^j\subset \Omega[X]\),
where \(W^j\) denotes the span of products. Multiplication of cusp
forms at these even weights identifies \(W^j\) with a subspace of
the generic evaluations of \(S_Z(ju)\). The Hilbert polynomial of the
projective image has leading term \(\delta_Z j^D/D!\). Dividing the
inclusion of dimensions by \(j^D/D!\) gives \(\delta_Z\le g_Z(u)\).

For the reverse comparison we spell out where the valuative estimate
enters. It supplies \(a\ge1\) and a multiplier \(\theta\) with nonzero
generic value. If \(k\ge1\), \(f\in S_Z(ak u')\), put
\(z=\operatorname{ev}(\theta^k f)\). The valuation inequalities on the
cone imply that \(zX^{ak}\) is integral over \(R\). Birationality and
the ratio condition give \(zX^{ak}\in\operatorname{Frac}(R)\). The latter
condition fixes the field in which normalization is taken. The integral
closure \(\overline R\) is finite over \(R\), since \(R\) is a finitely
generated domain over a number field. Clearing denominators of a finite
set of module generators gives a nonzero \(c\in R\) with
\(c\overline R\subset R\).

Choose a nonzero homogeneous coefficient \(c_e\) of \(c\), so that
\(c_e\in W^e\). The degree \(ak+e\) coefficient of \(czX^{ak}\in R\)
gives \(c_ez\in W^{ak+e}\). Multiplying by \(c_e^{a-1}\) yields
\(c_e^az\in W^{a(k+e)}\). All the multipliers have nonzero generic
value. Evaluation of a cusp form on the generic component is injective
on the space defined on that component. Consequently
\[
 \dim S_Z(ak u')\le \dim W^{a(k+e)} .
\]
Divide by \(k^D/D!\) and pass to the limit. By
\eqref{eq:volume-nef-growth} the left side tends, after normalization,
to \(a^Dg_Z(u')\). The Hilbert polynomial gives the limit
\(a^D\delta_Z\) on the right, since \(e\) is fixed. Cancelling \(a^D\)
proves the other half of \eqref{eq:volume-sandwich}. Equality of
component covolumes now allows \(Z'\) to replace \(Z\) in the growth
factors. In the window \(1-\rho=V(u')/V(u)>0\), so
\(g_Z(u')=(1-\rho)g_Z(u)\) by \eqref{eq:volume-flag}. This gives
\eqref{eq:rankshare-vol}.

\paragraph{Source correspondences and geometric inputs.}
The weight decomposition is \srcid{EqLastAddGaps}; nefness is
\srcid{NefLineBundleWeightGap}. The limit on the boundary and the
normalization of the covolume appear in \srcid{ExistsEffCartierCompl}
and \srcid{FlagDegreeFormula}. The conductor calculation is in
\srcid{ConeDegreeGrowthChowFormIntegrality}, and its application to cusp
forms is \srcid{CuspFormGrowthLeImageDegreeLowerWeight}. The two-sided
comparison is \srcid{ComponentImageDegreeBetweenOfProportionality}.
These arguments use the toroidal model and comparison theorem, nef
asymptotic Riemann--Roch, the comparison of self-dual lattices at good
level, and the existence of the valuative multiplier at \(u,u'\).
